%% kapa: exam=1 round=22 year=2011 period=1 form=A subject=economics subject_ko=경제원론 seq=32 number=72
%% source: Q-Net 감정평가사 기출문제 (cst00309) artlSeq=5036790 file=extracted/1008/5036790/2011년 제22회 감정평가사 1차시험 기출문제/1교시.hwp
%% conversion: kordoc 4.24.51 -> tools/gen1.py (auto); review=proofread-sonnet; flags=none
\begin{question}{72}
다음 ( )안의 용어가 순서대로 올바른 것은?
\begin{framedtext}
후방굴절하는 노동공급곡선은 여가-소득 선택모형에서 임금율의 변화에 따라 도출되는 ( ㄱ )소비곡선에서 유도되고, 소득효과와 대체효과를 비교할 경우 노동공급곡선의 우하향하는 구간에서는 ( ㄴ )효과가 더 크다.
\end{framedtext}
\begin{choices}
  \item ㄱ. 임금, ㄴ. 대체
  \item ㄱ. 가격, ㄴ. 소득
  \item ㄱ. 가격, ㄴ. 대체
  \item ㄱ. 소득, ㄴ. 소득
  \item ㄱ. 소득, ㄴ. 대체
\end{choices}
\end{question}
\answer{2}
